\section{Interiors of actual fan triangles}
\label{sec:area_law_fan_regularity}

Fix an arbitrary origin $o\in\mathbb R^2$, natural dyadic exponent $\ell$
and integer cell index $z\in\mathbb Z^2$. Choose any optional midpoint
subdivisions of the four sides. Write $c$ for the cell center and
$P_e=\operatorname{conv}\{c,a_e,b_e\}$ for each actual fan triangle;
see \cite[Section~11]{OpenAI2026AreaLaw}.

% Source: 10-geometry.tex, prop:two-families, lines 299--323,
% especially 308--316, and geometry:initial-stars, lines 352--363;
% revision adc7f1241b42e322a6451854ab7e4b4c146bf78a.

\begin{theorem}[Nonempty interiors and closure of fan triangles]
    \label{thm:area_law_cell_fan_regularity}
    \lean{TNLean.PEPS.AreaLaw.Geometry.cellFanPolygon_interior_nonempty_and_closure_eq}
    \leanok
    \uses{def:area_law_cell_fan}
    Every actual fan triangle satisfies
    \begin{align}
        \operatorname{int}(P_e)            & \ne\varnothing, \\
        \overline{\operatorname{int}(P_e)} & =P_e.
    \end{align}
\end{theorem}
\begin{proof}\leanok
    \uses{def:area_law_cell_fan,def:area_law_template}
    The defining determinant $\det(a_e-c,b_e-c)$ is nonzero. Thus the
    two edge vectors are linearly independent and span the plane, and
    the three vertices affinely span $\mathbb R^2$. Their convex hull
    consequently has nonempty interior.

    The convex hull of these finitely many vertices is closed. A convex
    set with nonempty interior has the same closure as its interior,
    so $\overline{\operatorname{int}(P_e)}=\overline{P_e}=P_e$.
\end{proof}
